% Standalone problem document, generated from the adjacent README.md.
% Compile with XeLaTeX. Mathematical content is maintained in Markdown.
\documentclass[11pt,a4paper]{article}
\usepackage[fontsize=10.5pt]{scrextend}
\usepackage[margin=22mm,headheight=15pt,headsep=9mm,footskip=11mm]{geometry}
\usepackage{fontspec}
\setmainfont{texgyrepagella}[Extension=.otf,UprightFont=*-regular,BoldFont=*-bold,ItalicFont=*-italic,BoldItalicFont=*-bolditalic]
\setsansfont{texgyreheros}[Extension=.otf,UprightFont=*-regular,BoldFont=*-bold,ItalicFont=*-italic,BoldItalicFont=*-bolditalic]
\setmonofont{texgyrecursor}[Extension=.otf,UprightFont=*-regular,BoldFont=*-bold,ItalicFont=*-italic,BoldItalicFont=*-bolditalic,Scale=MatchLowercase]
\usepackage{amsmath,amssymb}
\usepackage{unicode-math}
\setmathfont{texgyrepagella-math.otf}
\usepackage{xcolor,microtype,enumitem,needspace,fancyhdr}
\usepackage{bookmark}
\definecolor{ink}{HTML}{17344B}
\definecolor{accent}{HTML}{197D80}
\definecolor{muted}{HTML}{596773}
\hypersetup{colorlinks=true,linkcolor=accent,urlcolor=accent,pdftitle={IE-26 - Sharp
stability bounds for perturbed Fourier
interpolation},pdfauthor={Open Problems in Numerical Linear Algebra}}
\urlstyle{same}
\setlength{\parindent}{0pt}
\setlength{\parskip}{5pt plus 1pt minus 1pt}
\linespread{1.04}
\setlength{\emergencystretch}{3em}
\widowpenalty=10000
\clubpenalty=10000
\displaywidowpenalty=10000
\allowdisplaybreaks[1]
\setlist{nosep,leftmargin=1.5em}
\providecommand{\tightlist}{\setlength{\itemsep}{0pt}\setlength{\parskip}{0pt}}
\providecommand{\pandocbounded}[1]{#1}
\pagestyle{fancy}
\fancyhf{}
\fancyhead[L]{\sffamily\footnotesize\color{muted}OPEN PROBLEMS IN NUMERICAL LINEAR ALGEBRA}
\fancyhead[R]{\sffamily\bfseries\footnotesize\color{accent}IE-26}
\fancyfoot[L]{\parbox[b]{0.9\textwidth}{\sffamily\fontsize{8}{10}\selectfont\color{muted}Editorial ratings; a bounded search cannot certify that a problem remains unsolved.\\Literature check: 2026-09-12}}
\fancyfoot[R]{\sffamily\footnotesize\color{muted}\thepage}
\renewcommand{\headrulewidth}{0.3pt}
\renewcommand{\headrule}{\hbox to\headwidth{\color{accent}\leaders\hrule height \headrulewidth\hfill}}
\makeatletter
\renewcommand{\section}{\Needspace{5\baselineskip}\@startsection{section}{1}{0pt}{12pt plus 2pt minus 2pt}{4pt}{\sffamily\bfseries\large\color{ink}}}
\renewcommand{\subsection}{\Needspace{4\baselineskip}\@startsection{subsection}{2}{0pt}{9pt plus 1pt minus 1pt}{3pt}{\sffamily\bfseries\normalsize\color{ink}}}
\makeatother
\setcounter{secnumdepth}{0}
\begin{document}
{\sffamily\small\color{muted}Linear systems and elimination\par}
\vspace{3pt}
{\raggedright\hyphenpenalty=10000\exhyphenpenalty=10000\sffamily\bfseries\fontsize{21}{24}\selectfont\color{ink}Sharp
stability bounds for perturbed Fourier interpolation\par}
\vspace{7pt}
{\sffamily\small\textbf{Difficulty:} challenging\quad\textbf{Importance:} interesting
to the community\par}
{\sffamily\small\bfseries\color{accent}Status: Solved\par}
\vspace{4pt}
{\color{accent}\hrule height 0.8pt}
\vspace{6pt}
\textbf{Topic:} Conditioning of Fourier matrices and interpolation
operators

\textbf{Rating rationale:} The conjectures ask for sharp dimension
dependence under a fixed relative perturbation of every sampling point.
They concern the sensitivity of Fourier coefficient recovery and
interpolation, with consequences for spectral computation and nonuniform
sampling.

\subsection{Resolution --- 12 September
2026}\label{resolution-12-september-2026}

\textbf{Solved affirmatively.}
\href{https://github.com/ajt60gaibb/OpenProblemsInNLA/blob/main/linear-systems-and-elimination/IE-26/solution.md}{Theorem
1 and Sections 2--7 of the complete proof} establish both original
bounds: the Lebesgue estimate with one absolute constant for every
\(0<\alpha<1/2\), and the normalized square Fourier inverse-norm
estimate without a logarithmic loss for each fixed \(1/4<\alpha<1/2\).
Every \(N\ge2\) and every allowed collection of shifts is included. The
second constant may depend on \(\alpha\); no endpoint assertion at
\(\alpha=1/4\) is made.

\textbf{Author:} George Stepaniants, Department of Computing and
Mathematical Sciences, California Institute of Technology, Pasadena,
California, USA.
\href{https://github.com/ajt60gaibb/OpenProblemsInNLA/blob/main/linear-systems-and-elimination/IE-26/solution.pdf}{Proof
PDF} ·
\href{https://github.com/ajt60gaibb/OpenProblemsInNLA/blob/main/linear-systems-and-elimination/IE-26/solution.tex}{Standalone
TeX}.

The full argument passed a separate
\href{https://github.com/ajt60gaibb/OpenProblemsInNLA/blob/main/references/stepaniants-ie26-2026-09-12/verification/independent-review/IE-26-independent-review.md}{independent
Codex-agent mathematical review}, bound to the exact frozen source. The
coordinating agent contributed before freeze and is recorded separately,
not as the independent reviewer. Substantial AI assistance is disclosed;
no external human peer review or formal verification is asserted.

Austin and Trefethen retain credit for the conjectures and prior
analysis, with the related thesis and later results cited below. The
periodic Hilbert-transform weak \((1,1)\) inequality is the external
analytic input; the proof supplies the remaining grid and matrix
estimates. The
\href{https://github.com/ajt60gaibb/OpenProblemsInNLA/blob/main/references/stepaniants-ie26-2026-09-12/README.md}{submission
and verification record} contains exact provenance and the bounded
public eligibility check. The original target, earlier evidence and
references remain below. The ratings above describe the former open
question.

\newpage

\subsection{Statement}\label{statement}

For each integer \(N\ge2\), put \(m=2N+1\) and \(h=2\pi/m\). Fix
\(0<\alpha<1/2\), choose arbitrary real numbers \(s_{-N},\ldots,s_N\)
with \(|s_k|\le\alpha\), and put \(x_k=(k+s_k)h\), interpreted modulo
\(2\pi\). Define the normalized Fourier matrix

\[F_N(s)=m^{-1/2}[e^{ijx_k}]_{k,j=-N}^{N}.\]

The nodes are distinct. For \(y\in\mathbb C^m\), let \(T_Ny\) be the
unique trigonometric polynomial of degree at most \(N\) with
\((T_Ny)(x_k)=y_k\). Set

\[\Lambda_N(s)=\sup_{\|y\|_\infty\le1}\|T_Ny\|_{L^\infty[-\pi,\pi]},\qquad
\Gamma_N(s)=\|F_N(s)^{-1}\|_2.\]

\textbf{Austin--Trefethen conjectures.} Establish the following two
uniform bounds:

\begin{enumerate}
\def\labelenumi{\arabic{enumi}.}
\item
  There is an absolute \(C>0\) such that, for every \(N\ge2\),
  \(0<\alpha<1/2\), and admissible \(s\),

  \[\Lambda_N(s)\le C\frac{N^{2\alpha}-1}{\alpha(1-2\alpha)}.\]
\item
  For every fixed \(1/4<\alpha<1/2\), there is \(C_\alpha>0\),
  independent of \(N\) and \(s\), such that

  \[\Gamma_N(s)\le C_\alpha N^{4\alpha-1}.\]
\end{enumerate}

These are kept together as one perturbed-grid stability target. The
first is the source's conjectured replacement in equation (12); the
second is its conjecture on p.2119. The normalized matrix makes the
discrete norm convention explicit, following equation (1.6) of
Chen--Lin--Zhang. The second statement does not assert an endpoint bound
at \(\alpha=1/4\). The first has the usual logarithmic limit as
\(\alpha\downarrow0\).

\subsection{Evidence and relation between the
bounds}\label{evidence-and-relation-between-the-bounds}

Austin's thesis, Conjecture 3.5, states the first bound with a universal
constant and predicts a matching worst-grid growth rate. Conjecture 3.10
concerns the second norm. This is finite-dimensional conditioning of an
interpolation map, rather than restricted isometry of column subsets of
an equispaced Fourier matrix as in
\href{https://github.com/ajt60gaibb/OpenProblemsInNLA/blob/main/frames-and-matrix-designs/FR-02/README.md}{FR-02}.

Chen--Lin--Zhang's August 2026 preprint gives

\[\Gamma_N(s)\le C_\alpha N^{4\alpha-1}\log N\]

for \(1/4<\alpha<1/2\), together with a family having a matching lower
power \(N^{4\alpha-1}\). Thus the logarithmic factor remains in the
checked upper bound. The discrete Kadec theorem of Yu--Townsend treats
\(\alpha<1/4\), outside the second displayed target; oversampling treats
a rectangular matrix instead of the square matrix here. Norm conversion
alone loses powers of \(N\) and does not identify the two conjectures.

\newpage
\subsection{References and status
check}\label{references-and-status-check}

\begin{itemize}
\tightlist
\item
  A. P. Austin and L. N. Trefethen, \emph{Trigonometric interpolation
  and quadrature in perturbed points}, SIAM Journal on Numerical
  Analysis 55(5) (2017), 2113--2122.
  \href{https://doi.org/10.1137/16M1107760}{DOI};
  \href{https://people.maths.ox.ac.uk/trefethen/perturbed.pdf}{author
  PDF}. The conjectured sharpening of (12) is stated on pp.2115--2116;
  the second-norm conjecture is on p.2119.
\item
  A. P. Austin, \emph{Some New Results on and Applications of
  Interpolation in Numerical Computation}, D.Phil. thesis, University of
  Oxford, 2016.
  \href{https://personal.math.vt.edu/apaustin/pubs/DPhilThesis.pdf}{Author
  PDF}. Chapter 3, Conjectures 3.5 (p.75) and 3.10 (p.82), supplies the
  original norm definitions and uniformity. The thesis's endpoint claim
  is not imported; the later 2017 paper's supercritical interval is
  followed.
\item
  L. Chen, R. Lin and H. Zhang, \emph{The smallest singular value of
  nonuniform Fourier matrices},
  \href{https://arxiv.org/abs/2608.21960v1}{arXiv:2608.21960v1}, 22
  August 2026. \href{https://arxiv.org/html/2608.21960v1}{Full text}.
  Introduction, Theorems 1.4--1.5 and equations (1.6)--(1.7), explicitly
  connect the new bounds to the second conjecture.
\item
  A. Yu and A. Townsend, \emph{On the stability of unevenly spaced
  samples for interpolation and quadrature}, BIT Numerical Mathematics
  63 (2023), article 23.
  \href{https://doi.org/10.1007/s10543-023-00965-z}{DOI}. Theorem 2.1
  and §6.2 concern the subcritical and oversampled regimes.
\end{itemize}

On 2026-09-11, checked the 2017 author paper, Austin's full thesis, the
current August 2026 arXiv paper, and targeted searches for
Austin--Trefethen conjectures, perturbed-grid Lebesgue constants, and
subsequent proofs/counterexamples. No proof of the first displayed bound
or removal of the logarithm in the second was located. The August 2026
result also addresses a smoother-function quadrature convergence
question; that different question is not counted here. This is a bounded
check, and the preprint's proofs have not been independently verified.
\end{document}
